
\section{The authentic source residual and the necessary nonzero-data gate}
\label{sec:source-startup-XXIX}

This section corrects the target of the unforcing programme.
The source moment identities (A.11) and (A.19) control leading-profile
matching, not the physical forcing after cutoff. In the original
forced Navier--Stokes source \cite{openai2026},
Proposition 10.1, equation (10.4), forms the compactly supported
velocity $u_B=\operatorname{curl}(c{\mathcal A})+c{\mathcal B}e_\theta$,
with a temporal cutoff giving $u_B(0)=0$. Equation (10.5)
then defines the actual force
\[
f_B=\partial_t u_B+(u_B\cdot\nabla)u_B-\nu\Delta u_B+\nabla p_B.
\]
Appendix A is not a certificate that this vector residual,
or its divergence-free Leray projection, is zero.

\begin{theorem}[Projected residual and compulsory unforced startup]
\label{thm:xxix-startup}
Let $\Omega=\mathbb R^3$ or $\mathbb T^3$, $\nu>0$, and
$u_B$ be any smooth divergence-free finite-energy field
on $[0,T)$ with $u_B(0)=0$. Let $w$ be smooth
and divergence-free and put $u=u_B+w$.
For an arbitrary smooth pressure $p$, define the residual
\[
R=\partial_tu+(u\cdot\nabla)u-\nu\Delta u+\nabla p,
\]
and assume $\mathbb P R\in L^1([0,t];L^2(\Omega))$
for each $t<T$. Then
\begin{equation}\label{eq:xxix-source-budget}
\boxed{\quad
\|u_B(t)\|_2\le
\|w(0)\|_2+
\int_0^t\|\mathbb P R(s)\|_2\,ds+
\|w(t)\|_2.\quad}
\end{equation}
In particular, if $w(0)=0$ and the completed field solves
the \emph{unforced} Navier--Stokes equations, i.e.
$\mathbb P R\equiv0$ with the pressure reconstructed
appropriately, then $u\equiv0$ and necessarily
\begin{equation}\label{eq:xxix-no-go}
\boxed{\qquad w=-u_B\qquad\text{on }[0,T).}
\end{equation}
Consequently a convergent retarded perturbation series
with zero initial correction cannot remove the entire
forced source while retaining any nonzero part of
the forced velocity. This does not rule out a nonzero
unforced initial datum.
\end{theorem}
\begin{proof}
Apply the orthogonal Leray projection to the residual
equation. Testing against $u$ and integrating by parts
over space yields
\[
\frac12\frac{d}{dt}\|u(t)\|_2^2+
\nu\|\nabla u(t)\|_2^2
=\langle\mathbb P R(t),u(t)\rangle.
\]
The usual regularization of $\|u(t)\|_2$ at its zeros
and Cauchy--Schwarz imply
\[
\|u(t)\|_2\le\|u(0)\|_2+
\int_0^t\|\mathbb P R(s)\|_2\,ds.
\]
Since $u(0)=w(0)$ and
$u_B(t)=u(t)-w(t)$, the triangle inequality gives
\eqref{eq:xxix-source-budget}.
If $w(0)=0$ and $\mathbb P R=0$, the energy identity
gives $\|u(t)\|_2=0$ for every $t<T$. This proves
\eqref{eq:xxix-no-go}.
\end{proof}

\begin{corollary}[A quantitative non-gauge forcing floor for the authentic forced velocity]
\label{cor:xxix-real-force}
For the nontrivial compact forced source solution $u_B$
of Theorem 1.1 and Proposition 10.1 in \cite{openai2026},
at any fixed $t_*\in(0,1)$ at which
$\|u_B(t_*)\|_2=m_*>0$, its actual
Section 10 residual satisfies
\begin{equation}\label{eq:xxix-real-source-floor}
\boxed{\quad
\int_0^{t_*}\|\mathbb P f_B(s)\|_2\,ds\ge m_*>0.
\quad}
\end{equation}
Thus the entire physical forcing cannot be a pressure
gradient. Pure pressure-equivalent seven-ray modifications
of Editions XXV--XXVIII cannot remove this projected
force while leaving $u_B$ unchanged.
\end{corollary}
\begin{proof}
Apply Theorem~\ref{thm:xxix-startup} with $u=u_B$,
$w=0$, and $R=f_B$. Existence of a fixed positive
$t_*$ follows from the source theorem's nontrivial
late-time velocity, although its numerical $L^2$
value is not supplied by the Appendix A moments.
\end{proof}

\begin{proposition}[Exact three-dimensional periodic shear negative control]
\label{prop:xxix-shear}
On the spatial torus of period $2\pi$ in every coordinate,
take $\nu>0$ and, for $0\le t\le t_*$,
\[
u_B(x,y,z,t)=(t\sin y,0,0),\quad p_B=0.
\]
It is divergence-free, its convection vanishes identically,
and the complete physical residual is
\[
f_B=(1+\nu t)\sin y\,e_x.
\]
Moreover
\[
\|u_B(t)\|_2=2\pi^{3/2}t,\quad
\int_0^t\|f_B(s)\|_2\,ds
=2\pi^{3/2}(t+\nu t^2/2),
\]
and the exact kinetic-energy equality is
\[
\|u_B(t)\|_2^2+
2\nu\int_0^t\|\nabla u_B(s)\|_2^2\,ds
=2\int_0^t\langle f_B(s),u_B(s)\rangle\,ds.
\]
The retarded correction $w=-u_B$ begins from zero and
makes the completed field the trivial unforced solution,
not a singular flow.
\end{proposition}
\begin{proof}
Differentiate the sine mode explicitly:
$\Delta u_B=-t\sin y\,e_x$,
$(u_B\cdot\nabla)u_B=0$. The spatial integral
$\int_{\mathbb T^3}\sin^2y\,dx\,dy\,dz=4\pi^3$
gives the displayed norms and equality.
\end{proof}

\paragraph{Corrected theorem dependency and formalization boundary.}
The next actual PDE input is the complete Section 10
cutoff residual $f_B$ and its Leray projection,
not only the A.11/A.19 moment-root data.
A proposed smooth unforced correction must start with
nonzero solenoidal initial velocity $w(0)$, must satisfy
the exact projected equation throughout the early-time
cutoff interval and terminal region, and must preserve
the intended singularity despite the energy budget.
The finite Lean component accompanying this edition
formalizes only scalar and normed-space inequalities;
it does not formalize the Navier--Stokes energy identity,
Leray projection, source expressions or unforced
singularity. The exact three-dimensional manufactured
shear and independent 50,000 rational budget fixtures
pass their numerical/symbolic audit.
Nothing in this section proves the Millennium problem.
